\section*{Bab \ref{ch:b1:counting} --- Pencacahan}

\begin{solution}{exo:b1:counting:1}
Tahap yang saling bebas dan kaidah hasil kali: ada $26^2 \times 10^3 \times 26^2
= 26^4 \times 1000 = 456\,976\,000$ pelat. Bila keempat hurufnya
berbeda dua-dua, tahap hurufnya membentuk susunan-$4$ atas
abjad: $26 \times 25 \times 24 \times 23 = 358\,800$ cara, jadi ada
$358\,800 \times 1000 = 358\,800\,000$ pelat.
\end{solution}

\begin{solution}{exo:b1:counting:2}
\textsc{sepatu} mempunyai $6$ huruf yang berbeda: $6! = 720$ anagram.
\textsc{ananas} mempunyai $6$ huruf dengan pengulangan ($3$ huruf a, $2$ huruf n,
$1$ huruf s): setiap anagram ditentukan oleh posisi huruf a
($\binom 63$ pilihan), lalu posisi huruf n di antara $3$ tempat sisanya
($\binom 32$), sedangkan huruf s menempati tempat terakhir:
$\binom{6}{3}\binom{3}{2} = 20 \times 3 = 60$ anagram
(setara dengan $6!/(3!\,2!\,1!) = 60$).
\end{solution}

\begin{solution}{exo:b1:counting:3}
Seluruhnya: $\binom{12}{4} = 495$. Tepat $2$ perempuan: pilih mereka
($\binom 72 = 21$) lalu $2$ laki-laki ($\binom 52 = 10$): ada $210$ panitia.
Sekurang-kurangnya satu laki-laki: komplemen dari ``tanpa laki-laki'',
$\binom{12}{4} - \binom{7}{4} = 495 - 35 = 460$.
\end{solution}

\begin{solution}{exo:b1:counting:4}
\emph{Bulan lahir:} kotaknya adalah $12$ bulan; $13$ orang ke dalam $12$
kotak memaksa dua orang berada pada kotak yang sama
(\cref{cor:b1:counting:pigeonhole}).

\emph{Bilangan bulat berurutan:} kotaknya adalah $n$ pasangan $\{1,2\},
\{3,4\}, \dots, \{2n-1, 2n\}$, yang mempartisi $\intint{1}{2n}$.
Memilih $n + 1$ bilangan bulat menempatkan dua di antaranya pada pasangan yang sama, dan kedua
unsur sebuah pasangan itu berurutan.
\end{solution}

\begin{solution}{exo:b1:counting:5}
Untuk $1 \leq k \leq n$,
\[
k \binom nk = k\,\frac{n!}{k!\,(n-k)!}
= n\,\frac{(n-1)!}{(k-1)!\,(n-k)!} = n \binom{n-1}{k-1}.
\]
Dengan menjumlahkan lalu mengindeks ulang memakai $j = k - 1$:
\[
\sum_{k=0}^{n} k \binom nk = n \sum_{j=0}^{n-1} \binom{n-1}{j}
= n\, 2^{n-1}
\]
menurut \cref{prop:b1:counting:identities}. (Cara lain: turunkan
$(1+x)^n = \sum_k \binom nk x^k$ lalu ambil $x = 1$.)
\end{solution}

\begin{solution}{exo:b1:counting:6}
\omterm{def:b1:logic:map}{Pemetaan} naik tegas $f \colon \intint{1}{k} \to \intint{1}{n}$
ditentukan oleh petanya, yaitu \omterm{def:b1:logic:sets}{himpunan} bagian beranggota $k$ dari $\intint{1}{n}$ (daftarkan
\omterm{def:b1:logic:sets}{himpunan} bagian itu dalam urutan naik); sebaliknya setiap \omterm{def:b1:logic:sets}{himpunan} bagian beranggota $k$ memberi tepat
satu \omterm{def:b1:logic:map}{pemetaan} semacam itu. Jadi ada $\binom nk$ \omterm{def:b1:logic:map}{pemetaan} naik tegas.

Jika $f$ hanya naik, tulis $g(i) = f(i) + i - 1$. Maka $g$
naik tegas (di antara dua argumen berurutan, $f$ bertambah $\geq 0$
sedangkan $i - 1$ bertambah $1$) dengan nilai di $\intint{1}{n + k - 1}$; dan
$f(i) = g(i) - i + 1$ memulihkan $f$ dari sebarang $g$ yang naik tegas
ke dalam $\intint{1}{n+k-1}$. Ini sebuah bijeksi, jadi ada
$\binom{n + k - 1}{k}$ \omterm{def:b1:logic:map}{pemetaan} naik.
\end{solution}

\begin{solution}{exo:b1:counting:7}
Belah \omterm{def:b1:logic:sets}{himpunan} $E$ beranggota $m + n$ unsur atas blok $M$ ($m$ unsur)
dan $N$ ($n$ unsur). \omterm{def:b1:logic:sets}{Himpunan} bagian beranggota $k$ dari $E$ memuat sejumlah $j$ unsur
$M$ ($0 \leq j \leq k$) dan $k - j$ unsur $N$; untuk $j$ yang tetap ada
$\binom mj \binom{n}{k-j}$ \omterm{def:b1:logic:sets}{himpunan} bagian semacam itu, dan kasus $j = 0, \dots,
k$ mempartisi \omterm{def:b1:logic:sets}{himpunan} bagian beranggota $k$ tadi. Kaidah jumlah memberikan
kesamaan Vandermonde\index{kesamaan Vandermonde}.

Dengan $m = n = k$: $\binom{2n}{n} = \sum_{j=0}^{n} \binom nj
\binom{n}{n-j} = \sum_{j=0}^{n} \binom nj^2$, memakai $\binom{n}{n-j} =
\binom nj$.
\end{solution}

\begin{solution}{exo:b1:counting:8}
Misalkan $A_d$ \omterm{def:b1:logic:sets}{himpunan} kelipatan $d$ di $\intint{1}{1000}$, sehingga
$\abs{A_d} = \lfloor 1000/d \rfloor$. Inklusi--eksklusi
(\cref{thm:b1:counting:inclexcl}) dengan $A_2, A_3, A_5$, sambil memperhatikan
$A_2 \cap A_3 = A_6$ dan seterusnya:
\[
500 + 333 + 200 - 166 - 100 - 66 + 33 = 734 .
\]
Jadi ada $734$ bilangan bulat yang habis dibagi $2$, $3$ atau $5$.
\end{solution}

\begin{solution}{exo:b1:counting:9}
Pada $2$ unsur: seluruh $2^4 = 16$ \omterm{def:b1:logic:map}{pemetaan} kecuali $2$ \omterm{def:b1:logic:map}{pemetaan} konstan:
ada $14$ surjeksi.

Pada $3$ unsur: dengan inklusi--eksklusi atas nilai yang terlewat,
banyaknya \omterm{def:b1:logic:map}{pemetaan} dari \omterm{def:b1:logic:sets}{himpunan} beranggota $4$ ke \omterm{def:b1:logic:sets}{himpunan} beranggota $3$ yang melewatkan sekurang-kurangnya satu nilai
adalah $\binom 31 2^4 - \binom 32 1^4 = 48 - 3 = 45$; seluruh \omterm{def:b1:logic:map}{pemetaan} berjumlah $3^4 =
81$; jadi surjeksinya: $81 - 45 = 36$. (Periksa: surjeksi dari $4$ pada
$3$ unsur menggandakan tepat satu nilai: pilih nilai yang digandakan
($3$), pasangan yang terpetakan ke sana ($\binom 42 = 6$), lalu bijeksi bagi
sisanya ($2$): $3 \times 6 \times 2 = 36$.)
\end{solution}

\begin{solution}{exo:b1:counting:10}
Sebuah penyelesaian $x_1 + \dots + x_n = k$ di $\N^n$ tersandikan sebagai sebaris
$k$ bintang dan $n - 1$ sekat: tulis $x_1$ bintang, sebuah sekat, $x_2$ bintang, sebuah
sekat, \dots, lalu diakhiri $x_n$ bintang. Ini bijeksi pada
kata sepanjang $k + n - 1$ yang memakai $k$ bintang dan $n - 1$ sekat, dan
kata semacam itu ditentukan oleh posisi bintangnya:
$\binom{n + k - 1}{k}$. Pemilihan dengan pengulangan berpadanan dengan
penyelesaian persamaan itu ($x_i$ = banyaknya salinan objek $i$), jadi
cacahannya sama.
\end{solution}

\begin{solution}{exo:b1:counting:11}
Dengan $A_i = \{\sigma : \sigma(i) = i\}$, sebuah \omterm{def:b1:counting:objects}{permutasi} di
$\bigcap_{i \in I} A_i$ menetapkan setiap $i \in I$ dan mempermutasikan
$n - \abs I$ titik lainnya secara bebas: $\abs{\bigcap_{i \in I} A_i} =
(n - \abs I)!$. Inklusi--eksklusi:
\[
\Bigl|\bigcup_i A_i\Bigr|
= \sum_{k=1}^{n} (-1)^{k+1} \binom nk (n-k)!
= \sum_{k=1}^{n} (-1)^{k+1} \frac{n!}{k!} ,
\]
karena ada $\binom nk$ \omterm{def:b1:logic:sets}{himpunan} bagian $I$ berukuran $k$. Jadi
\[
D_n = n! - \Bigl|\bigcup_i A_i\Bigr|
= n!\Bigl(1 - \sum_{k=1}^{n} \frac{(-1)^{k+1}}{k!}\Bigr)
= n! \sum_{k=0}^{n} \frac{(-1)^k}{k!} .
\]
Untuk kesamaan kedua: golongkan \omterm{def:b1:counting:objects}{permutasi} $\sigma$ dari
$\intint{1}{n}$ menurut \omterm{def:b1:logic:sets}{himpunan} titik tetapnya $F(\sigma)$. Untuk
\omterm{def:b1:logic:sets}{himpunan} bagian beranggota $k$ yang tetap, katakanlah $F$, \omterm{def:b1:counting:objects}{permutasi} dengan $F(\sigma) = F$ tepat adalah
\omterm{def:b1:counting:objects}{permutasi} kacau pada komplemennya: ada $D_{n-k}$ buah. Menjumlahkan atas
$\binom nk$ pilihan $F$ untuk setiap $k$ memberikan
$n! = \sum_{k=0}^{n} \binom nk D_{n-k}$.
\end{solution}

\begin{solution}{exo:b1:counting:12}
\emph{Kesamaan pertama.} Cacah pasangan $(A, a)$ dengan $A \subseteq E$
($\abs E = n$) dan $a \in A$. Menurut ukuran $A$: ada $\sum_k \binom nk k$
pasangan. Dengan memilih unsur bertanda lebih dulu: ada $n$ pilihan bagi $a$, lalu
sebarang \omterm{def:b1:logic:sets}{himpunan} bagian dari $n - 1$ unsur sisanya untuk melengkapi $A$:
ada $n\,2^{n-1}$ pasangan.

\emph{Kesamaan kedua.} Cacah tripel $(A, a, b)$ dengan $a, b \in
A$ (boleh $a = b$). Menurut ukurannya: $\sum_k k^2 \binom nk$. Secara langsung:
entah $a = b$ (ada $n\,2^{n-1}$ tripel, cacahan sebelumnya) atau $a \neq b$
(ada $n(n-1)$ pilihan terurut, lalu sebarang \omterm{def:b1:logic:sets}{himpunan} bagian dari $n - 2$ unsur
lainnya: $n(n-1)\,2^{n-2}$). Totalnya
\[
n\,2^{n-1} + n(n-1)\,2^{n-2} = n\,2^{n-2}\,(2 + n - 1)
= n(n+1)\,2^{n-2} .
\]
\end{solution}

\begin{solution}{pb:b1:counting:1}
\textbf{1.} $D_1 = 0$ (satu-satunya \omterm{def:b1:counting:objects}{permutasi} menetapkan $1$), $D_2 = 1$
(pertukarannya), $D_3 = 2$ (dalam notasi satu baris: $231$ dan $312$). Untuk
$n = 4$, kelompokkan menurut $\sigma(1)$: dengan $\sigma(1) = 2$ \omterm{def:b1:counting:objects}{permutasi} kacaunya
adalah $2143$, $2341$, $2413$; dengan $\sigma(1) = 3$: $3142$, $3412$,
$3421$; dengan $\sigma(1) = 4$: $4123$, $4312$, $4321$. Tiga buah pada masing-masing
kelompok: $D_4 = 9$.

\textbf{2.} Sebuah \omterm{def:b1:counting:objects}{permutasi} dengan tepat $k$ titik tetap
ditentukan oleh pilihan \omterm{def:b1:logic:sets}{himpunan} titik tetapnya $F$ ($\binom nk$
cara) beserta pembatasannya pada komplemen, yang haruslah
\omterm{def:b1:counting:objects}{permutasi} $n - k$ titik \emph{tanpa} titik tetap
($D_{n-k}$ cara). Kedua pilihan itu saling bebas dan
padanannya \omterm{def:b1:logic:inj}{bijektif}: $P_k(n) = \binom nk D_{n-k}$.

\textbf{3.} $P_0(4) = D_4 = 9$; $P_1(4) = \binom41 D_3 = 4 \times 2
= 8$; $P_2(4) = \binom42 D_2 = 6$; $P_3(4) = \binom43 D_1 = 0$
(tiga titik tetap memaksa adanya titik tetap yang keempat); $P_4(4) = 1$. Jumlahnya: $9 + 8 + 6
+ 0 + 1 = 24 = 4!$. Tak ada yang cocok ($9$ kasus) mengungguli tepat satu yang cocok
($8$ kasus) --- tipis saja.

\textbf{4.} Cacah pasangan $(\sigma, i)$ dengan $\sigma(i) = i$. Untuk
$i$ yang tetap, \omterm{def:b1:counting:objects}{permutasi} yang menetapkan $i$ adalah \omterm{def:b1:counting:objects}{permutasi}
$n - 1$ titik lainnya: ada $(n-1)!$ buah. Jadi banyaknya pasangan
adalah $n \cdot (n-1)! = n!$, dan bilangan itu juga sama dengan
$\sum_\sigma \abs{\mathrm{Fix}(\sigma)}$. Dengan membaginya dengan banyaknya
\omterm{def:b1:counting:objects}{permutasi}, yaitu $n!$: rata-rata banyaknya titik tetap adalah
tepat $1$, untuk setiap $n \geq 1$.

\textbf{5.} Misalkan $\sigma$ \omterm{def:b1:counting:objects}{permutasi} kacau $\intint1{n+1}$ dan
$j = \sigma(n+1) \in \intint1n$: ada $n$ nilai yang mungkin. \emph{Kasus
$\sigma(j) = n+1$:} titik $j$ dan $n+1$ bertukar, sedangkan $\sigma$
yang dibatasi pada $n - 1$ titik sisanya adalah \omterm{def:b1:counting:objects}{permutasi} kacau
sembarang atas titik-titik itu: ada $D_{n-1}$ kemungkinan. \emph{Kasus $\sigma(j)
\neq n+1$:} tulis $i_0 = \sigma^{-1}(n+1)$; di sini $i_0 \neq j$ dan
$i_0 \leq n$. Definisikan $\tau$ pada $\intint1n$ dengan $\tau(i) = \sigma(i)$
untuk $i \neq i_0$ dan $\tau(i_0) = j$. Maka $\tau$ adalah \omterm{def:b1:counting:objects}{permutasi}
$\intint1n$ (nilai $n+1$ telah digantikan oleh nilai
$j$ yang hilang), dan ia \omterm{def:b1:counting:objects}{permutasi} kacau: $\tau(i_0) = j \neq i_0$, dan
$\tau(i) = \sigma(i) \neq i$ di tempat lain. Sebaliknya, dari sebuah
\omterm{def:b1:counting:objects}{permutasi} kacau $\tau$ atas $\intint1n$ dan nilai $j$, kita pulihkan
$\sigma$ dengan menetapkan $\sigma(n+1) = j$, $\sigma(\tau^{-1}(j)) = n+1$
dan $\sigma = \tau$ di tempat lain: sebuah bijeksi, yang memberikan $D_n$
kemungkinan. Dengan menjumlahkan atas $j$: $D_{n+1} = n(D_n + D_{n-1})$.
Secara numerik: $D_5 = 4(9 + 2) = 44$, $D_6 = 5(44 + 9) = 265$.

\textbf{6.} Dari pertanyaan 5, $D_{n+1} = nD_n + nD_{n-1}$, jadi
\[
u_{n+1} = D_{n+1} - (n+1)D_n = nD_n + nD_{n-1} - (n+1)D_n
= -(D_n - nD_{n-1}) = -u_n .
\]
Karena $u_1 = D_1 - 1 \cdot D_0 = -1$, induksi memberikan $u_n =
(-1)^n$, yakni $D_n = nD_{n-1} + (-1)^n$ untuk $n \geq 1$.

\textbf{7.} Induksi pada $n$. Basis: $D_0 = 1 = 0!\,s_0$. Langkah:
dengan mengandaikan $D_{n-1} = (n-1)!\,s_{n-1}$,
\[
D_n = nD_{n-1} + (-1)^n = n!\,s_{n-1} + (-1)^n
= n!\Bigl(s_{n-1} + \frac{(-1)^n}{n!}\Bigr) = n!\,s_n ,
\]
yang tepat merupakan rumus itu. Tak ada inklusi--eksklusi yang dipakai: hanya
rekursi kombinatorial pertanyaan 5.

\textbf{8.} Revisi trinomial, lewat faktorial:
\[
\binom nk \binom kj
= \frac{n!}{k!\,(n-k)!} \cdot \frac{k!}{j!\,(k-j)!}
= \frac{n!}{j!\,(n-j)!} \cdot \frac{(n-j)!}{(k-j)!\,(n-k)!}
= \binom nj \binom{n-j}{k-j} .
\]
Sekarang substitusikan $a_k = \sum_j \binom kj b_j$ lalu tukarkan kedua
jumlah \omterm{def:b1:counting:card}{hingga} itu:
\[
\sum_{k=0}^{n} (-1)^{n-k} \binom nk a_k
= \sum_{j=0}^{n} b_j \binom nj
\sum_{k=j}^{n} (-1)^{n-k} \binom{n-j}{k-j}
= \sum_{j=0}^{n} b_j \binom nj
\sum_{i=0}^{n-j} (-1)^{(n-j)-i} \binom{n-j}{i} .
\]
Jumlah bagian dalamnya adalah penjabaran $(1 + (-1))^{n-j} = 0^{n-j}$
(teorema binomial, \cref{thm:b1:counting:binomial}): ia lenyap
untuk $j < n$ dan bernilai $1$ untuk $j = n$. Hanya $j = n$ yang bertahan, dan
ruas kanannya adalah $b_n$, sesuai yang diklaim.

\textbf{9.} Menurut simetri $\binom nk = \binom n{n-k}$, kesamaan
pada \cref{exo:b1:counting:11} dapat ditulis ulang sebagai $n! =
\sum_{k=0}^n \binom nk D_k$. Terapkan pertanyaan 8 dengan $a_n = n!$ dan
$b_k = D_k$:
\[
D_n = \sum_{k=0}^{n} (-1)^{n-k} \binom nk k!
= \sum_{k=0}^{n} (-1)^{n-k} \frac{n!}{(n-k)!}
= n! \sum_{j=0}^{n} \frac{(-1)^j}{j!} ,
\]
dengan mengindeks ulang lewat $j = n - k$: rumus itu untuk ketiga kalinya.

\textbf{10.} $D_n = n!\,s_n$ (pertanyaan 7), jadi
\[
\Bigl| D_n - \frac{n!}{\eu} \Bigr|
= n!\,\abs{s_n - \eu^{-1}} < \frac{n!}{(n+1)!} = \frac1{n+1} .
\]

\textbf{11.} Untuk $n \geq 1$ berlaku $\frac1{n+1} \leq \frac12$, dan
ketaksamaan pertanyaan 10 bersifat tegas: $D_n$ terletak pada jarak
$< \frac12$ dari $n!/\eu$, jadi ia satu-satunya bilangan bulat terdekat. Untuk
$n = 0$ batas itu hanya memberikan jarak $< 1$, dan klaimnya memang
gagal di situ: $0!/\eu \approx 0.368$ mempunyai bilangan bulat terdekat $0$, sedangkan
$D_0 = 1$.

\textbf{12.} $\eu^{-1} - s_n = \sum_{k \geq n+1} (-1)^k/k!$ adalah
deret berselang-seling dengan suku yang menurun tegas, jadi tandanya sama dengan
tanda suku pertamanya $(-1)^{n+1}/(n+1)!$. Jadi $s_n -
\eu^{-1}$ bertanda $(-1)^n$: untuk $n$ genap, $s_n > \eu^{-1}$
dan $D_n = n!\,s_n > n!/\eu$; untuk $n$ ganjil, $D_n < n!/\eu$.

\textbf{13.} $D_7 = 6(265 + 44) = 6 \times 309 = 1854$; $D_8 =
7(1854 + 265) = 7 \times 2119 = 14\,833$; $D_9 = 8(14\,833 + 1854)
= 8 \times 16\,687 = 133\,496$; $D_{10} = 9(133\,496 + 14\,833) =
9 \times 148\,329 = 1\,334\,961$. Periksa: $10!/\eu = 3\,628\,800 /
2.718281828 \approx 1\,334\,960.92$, yang bilangan bulat terdekatnya adalah
$1\,334\,961$ --- dan $D_{10} > 10!/\eu$, sebagaimana diramalkan pertanyaan 12
untuk $n$ genap.

\textbf{14.} $\abs{p_n - \eu^{-1}} = \abs{s_n - \eu^{-1}} <
\frac1{(n+1)!}$. Untuk $n = 6$: $p_6 = 265/720 = 0.36806$ (lima
angka desimal), berbanding $\eu^{-1} = 0.36788$; selisihnya di bawah $1/7! =
1/5040 < 2 \times 10^{-4}$. Batas $1/(n+1)!$ runtuh demikian cepat
sehingga peluangnya sudah terpaku sampai banyak angka desimal bahkan untuk
selusin surat: jawaban ``sekitar $36.8\%$'' itu, untuk setiap
keperluan praktis, tak bergantung pada $n$ --- kejutan yang termasyhur
dari soal ini.

\textbf{15.} Menurut pertanyaan 2 dan $D_m = m!\,s_m$:
\[
\frac{P_k(n)}{n!} = \frac{\binom nk D_{n-k}}{n!}
= \frac{D_{n-k}}{k!\,(n-k)!} = \frac{s_{n-k}}{k!}
\;\longrightarrow\; \frac{\eu^{-1}}{k!}
\]
ketika $n \to \infty$ dengan $k$ yang tetap, karena $s_{n-k} \to \eu^{-1}$.
Nilai limit $\eu^{-1}/k!$ ($k \in \N$) adalah bobot
distribusi Poisson berparameter $1$.

\textbf{16.} Cacah tripel $(\sigma, i, j)$ dengan $i \neq j$,
$\sigma(i) = i$, $\sigma(j) = j$. Dengan memilih pasangan terurutnya lebih dulu:
ada $n(n-1)$ cara; \omterm{def:b1:counting:objects}{permutasi} yang menetapkan $i$ sekaligus $j$ adalah
\omterm{def:b1:counting:objects}{permutasi} $n - 2$ titik sisanya: ada $(n-2)!$ buah.
Totalnya: $n(n-1)(n-2)! = n!$. Sebaliknya, menjumlahkan atas $\sigma$ lebih dulu
mencacah, untuk setiap $\sigma$, pasangan terurut titik tetap yang
berbeda: $\abs{\mathrm{Fix}(\sigma)}(\abs{\mathrm{Fix}(\sigma)}-1)$.
Jadi kesamaan yang dinyatakan itu berlaku; dengan membaginya dengan $n!$, rata-rata
$\abs{\mathrm{Fix}}(\abs{\mathrm{Fix}} - 1)$ adalah $1$, sehingga rata-rata
$\abs{\mathrm{Fix}}^2$ adalah $1 + 1 = 2$ dan ragamnya $2 -
1^2 = 1$.

\textbf{17.} Proporsi $1 - p_n$: untuk $n = 4$, $1 - \frac
9{24} = \frac{15}{24} = 0.6250$; untuk $n = 5$, $1 - \frac{44}{120} =
\frac{76}{120} = 0.6333$; untuk $n = 6$, $1 - \frac{265}{720} =
\frac{455}{720} = 0.6319$. Semuanya dalam jarak satu persen dari $1 - \eu^{-1}
\approx 0.6321$, dan berayun di sekitarnya.

\textbf{18.} Secara langsung:
\[
s_{n+2} - s_n = \frac{(-1)^{n+1}}{(n+1)!} +
\frac{(-1)^{n+2}}{(n+2)!}
= (-1)^{n+1}\Bigl(\frac1{(n+1)!} - \frac1{(n+2)!}\Bigr),
\]
sedangkan kurungnya bernilai $> 0$. Untuk $n$ genap selisihnya
negatif: $s_{n+2} < s_n$, jadi $p_0 > p_2 > p_4 > \dots$; untuk $n$
ganjil ia positif: $p_1 < p_3 < p_5 < \dots$ Digabungkan dengan
pertanyaan 12 (yang genap di atas $\eu^{-1}$, yang ganjil di bawah) dan pertanyaan 14
(jarak ke $\eu^{-1}$ menuju $0$): kedua tangga itu menjepit
$\eu^{-1}$ di antara keduanya.

\textbf{19.} Satu penarikan lengkap adalah \omterm{def:b1:counting:objects}{permutasi} acak seragam,
yang sah bila ia \omterm{def:b1:counting:objects}{permutasi} kacau: peluangnya $p_n \approx \eu^{-1}$.
Menurut fakta yang dikutip tadi, rata-rata banyaknya penarikan sampai berhasil adalah
$1/p_n$, dan pertanyaan 14 memberikan $1/p_n \approx \eu$ dengan galat
yang sudah dapat diabaikan bahkan untuk $n$ yang kecil. Jadi tukar kado rahasia dengan
pengulangan menelan rata-rata sekitar $\eu \approx 2.72$ penarikan
lengkap --- entah kantornya berisi $6$ orang atau $600$.

\textbf{20.} Tetapkan $j \geq 2$ lalu jalankan penggolongan pertanyaan 5 pada
nilai $\sigma(1) = j$. Jika $\sigma(j) = 1$: sisa
$n - 2$ titik memikul \omterm{def:b1:counting:objects}{permutasi} kacau sembarang, ada $D_{n-2}$ cara. Jika
$\sigma(j) \neq 1$: alihkan \omterm{def:b1:logic:map}{prapeta} $i_0 = \sigma^{-1}(1)$ ke
$j$ persis seperti pada pertanyaan 5; ini bijeksi dengan
\omterm{def:b1:counting:objects}{permutasi} kacau atas $n - 1$ titik $\{2, \dots, n\}$: ada $D_{n-1}$
cara. Totalnya $D_{n-1} + D_{n-2}$, sama untuk setiap $j$. Dengan menjumlahkan
atas $n - 1$ nilai $j$: $D_n = (n-1)(D_{n-1} + D_{n-2})$,
yang memperlihatkan faktor $n - 1$: $(n-1) \mid D_n$.

\textbf{21.} Klaimnya: $D_n$ ganjil jika dan hanya jika $n$ genap. Induksi memakai
$D_n = nD_{n-1} + (-1)^n$, yakni $D_n \equiv nD_{n-1} + 1 \pmod 2$.
Basis: $D_1 = 0$ genap sedangkan $n = 1$ ganjil: klaimnya berlaku. Jika $n$ genap,
$nD_{n-1}$ genap dan $D_n \equiv 1$: ganjil, sesuai klaim. Jika $n$
ganjil, maka $n - 1$ genap, sehingga $D_{n-1}$ ganjil menurut hipotesis, dan
$D_n \equiv D_{n-1} + 1 \equiv 0$: genap. Induksinya pun tertutup.

\textbf{22.} Mereduksi $D_n = nD_{n-1} + (-1)^n$ modulo $n$ mematikan
suku pertamanya: $D_n \equiv (-1)^n \pmod n$. Untuk $n = 10$:
$(-1)^{10} = 1$, dan memang $D_{10} = 1\,334\,961$ berakhir pada
angka $1$.

\textbf{23.} Untuk $n \geq 3$ berlaku $D_{n-1} \geq 1$, dan membagi
rekursi pertanyaan 6 dengan $D_{n-1}$ memberikan $D_n/D_{n-1} = n +
(-1)^n/D_{n-1}$, dengan $\abs{(-1)^n/D_{n-1}} \leq 1$ dan menuju
$0$ dengan cepat. Kesejalanannya: jika $D_n \approx n!/\eu$, maka
$D_n/D_{n-1} \approx n!/(n-1)! = n$ --- faktor $\eu$ saling hapus
pada rasionya, dan rekursi itu membenarkannya sampai ketelitian
$1/D_{n-1}$.

\textbf{24.} (i) Kaidah hasil kali dan kaidah jumlah melandasi setiap cacahan:
pertanyaan 2 dan 5 mempartisi \omterm{def:b1:logic:sets}{himpunan} \omterm{def:b1:counting:objects}{permutasi} atas tahap-tahap yang
saling bebas. (ii) Pencacahan ganda memberikan rata-rata (pertanyaan 4) dan
ragam (pertanyaan 16) banyaknya titik tetap tanpa
rumus $D_n$ sama sekali. (iii) Teorema binomial menghitung
jumlah dalam yang berselang-seling $(1-1)^{n-j}$, yang membuat inversi binomial
berjalan (pertanyaan 8). (iv) Batas deret berselang-seling mengubah
jumlah $n!\,s_n$ yang eksak tetapi buram menjadi \omterm{def:b1:logic:statement}{pernyataan} yang bening,
yaitu ``bilangan bulat terdekat dengan $n!/\eu$'' (pertanyaan 10--14).

\textbf{25.} Inklusi--eksklusi
(\cref{ex:b1:counting:derangement} dan \cref{exo:b1:counting:11})
adalah bukti yang konseptual: ia menjelaskan jumlah berselang-seling itu sebagai
koreksi atas pencacahan yang berlebih, dan ia dapat diperluas kata demi kata ke pencacahan
unsur yang menghindari sebarang keluarga \omterm{def:b1:logic:sets}{himpunan} ``buruk''. Jalur rekursi
(pertanyaan 5--7) menghitung paling cepat --- waktu linear, aritmetika
bilangan bulat yang eksak, tanpa faktorial --- dan menjadi sumber fakta
aritmetika pada Bagian V. Inversi binomial (pertanyaan 8--9) menempatkan
rumus itu di dalam transformasi umum yang akan muncul kembali di mana pun
dua sistem kesamaan segitiga saling berhadapan. Dan
munculnya $\eu$ paling baik dijelaskan oleh rumus itu sendiri: proporsi
\omterm{def:b1:counting:objects}{permutasi} kacau adalah jumlah parsial $s_n$ dari deret
untuk $\eu^{-1}$, jadi amplop Montmort itu, tiga dasawarsa sebelum
notasi Euler, sudah menghitung bilangan $\eu$.

\end{solution}
